% Intended LaTeX compiler: pdflatex
\documentclass[a4paper,12pt]{article}
\usepackage[utf8]{inputenc}
\usepackage[T1]{fontenc}
\usepackage{graphicx}
\usepackage{longtable}
\usepackage{wrapfig}
\usepackage{rotating}
\usepackage[normalem]{ulem}
\usepackage{amsmath}
\usepackage{amssymb}
\usepackage{capt-of}
\usepackage{hyperref}
\usepackage{\string~/.emacs.d/common}
\author{mahmood}
\date{\today}
\title{probability}
\hypersetup{
 pdfauthor={mahmood},
 pdftitle={probability},
 pdfkeywords={},
 pdfsubject={probability theory},
 pdfcreator={Emacs 28.2 (Org mode 9.5.5)}, 
 pdflang={English}}
\begin{document}

\maketitle
\begin{definition}
\textbf{probability} is a \href{discrete_maths2.org}{function}, denoted by \(P\) that describes the probability for a specific incident to occur\\
\begin{lemma}
the probability of an impossible event is zero; the probability of a certain event is one. therefore, for any event \(A\), the range of possible probabilities is \(0 \le P(A) \le 1\)\\
\end{lemma}
\begin{lemma}
for \(S\) the sample space of all possibilities, \(P(S)=1\). that is, the sum of all the probabilities for all possible events is equal to one.\\
\begin{entailment}
the sum of the probabilities of both \(A\) and \textbf{not} \(A\) is equal to 1\\
\end{entailment}
\end{lemma}
\begin{lemma}
\[
P(\varnothing) = 0
\]\\
\end{lemma}
\begin{lemma}
\[
P(\overline{A}) = 1 - P(A)
\]\\
\end{lemma}
\begin{lemma}
\[
A \subset B \implies P(A) < P(B)
\]\\
\end{lemma}
\begin{lemma}
let \(\Omega\) be a universe that contains \(n\) possibilities, and let \(A \subseteq \Omega\) such that \(A\) contains m elements, then the possibility for \(A\) to occur is \(P(A) = \frac{m}{n}\)\\
\begin{my_example}
we pick a random number between 0 and 9 (including both)\\
which means that \(\Omega=\{0,1,2,3,4,5,6,7,8,9\}\)\\
since there are 10 numbers, the odds of picking a number bigger than 5, which could be one of \(A=\{6,7,8,9\}\) are \(\frac{4}{10}\) because here \(A\) contains 4 elements and \(\Omega\) contains 10 elements\\
\end{my_example}
\end{lemma}
\begin{my_example}
in a family there are 3 kids\\
we define the following expressions:\\
E - the oldest is a boy\\
F - the younger kid (not youngest) is a boy\\
G - the youngest is a boy\\
find expressions for the following cases:\\
\begin{subquestion}
the three kids are boys\\
\begin{answer}
\[
E \cap F \cap G
\]\\
\end{answer}
\end{subquestion}
\begin{subquestion}
the three kids are girls\\
\begin{answer}
\[
\overline{E} \cap \overline{F} \cap \overline{G}
\]\\
\end{answer}
\end{subquestion}
\begin{subquestion}
atleast one of the kids is a boy\\
\begin{answer}
\[
E \cup F \cup G
\]\\
\end{answer}
\end{subquestion}
\begin{subquestion}
the two oldest kids are boys\\
\begin{answer}
\[
E \cap F
\]\\
\end{answer}
\end{subquestion}
\begin{subquestion}
exactly two kids are boys\\
\begin{answer}
\[
(E \cap F \cap \overline{G}) \cup (E \cap \overline{F} \cap G) \cup (\overline{E} \cap F \cap G)
\]\\
\end{answer}
\end{subquestion}
\end{my_example}
\begin{question}
given \(A \cup B\) is a union of 2 independent events, show that \(P(A \cup B) = P(A) + P(B) - P(A \cap B)\)\\
\begin{answer}
could use some venn diagrams here..\\
\begin{gather*}
  A \cup B = A \cup (B - A) \implies P(A \cup B) = P(A \cup (B - A)) = P(A) + P(B - A)\\
  B = (A \cap B) \cup (B - A) \implies P(B) = P(A \cap B) + P(B - A) \implies P(B - A) = P(B) - P(A \cap B)\\
  P(A \cup B) = P(A) + P(B) - P(A \cap B)
\end{gather*}
\end{answer}
\end{question}
\begin{question}
prove:\\
\[
A \subset B \implies P(B-A) = P(B)-P(A)
\]\\
\begin{answer}
\begin{gather*}
  B = (B \cap A) \cup (B - A) \implies P(B) = P(B \cap A) + P(B - A)\\
  A \subset B \implies A \cap B = A \implies P(B) = P(A) + P(B - A) \implies P(B - A) - P(B) - P(A)
\end{gather*}
\end{answer}
\end{question}
\begin{question}
let \(\Omega = \{1,2,3,4,5,6\}\) such that \(n=1,2,3,4,5,6\) and \(P(\{n\}) = \alpha n\), find \(\alpha\)\\
\begin{answer}
\[
P(\Omega) = 1 \implies \alpha + 2\alpha + \dots + 6\alpha = 1 \implies 21\alpha = 1 \implies \alpha = \frac{1}{21}
\]\\
\end{answer}
\end{question}
\end{definition}
\section{probability space}
\label{sec:org75774fa}
\begin{definition}
a \textbf{probability space} \((\Omega ,{\mathcal {F}},P)\) is a mathematical construct that provides a formal model of a random process or "experiment".\\
a probability space consists of three elements:\\
\begin{enumerate}
\item a \textbf{sample space}, \(\Omega\), which is the \href{discrete_maths2.org}{set} of all possible outcomes\\
\item an \textbf{event space}, which is a set of events \(\mathcal{F}\), an event being a set of outcomes in the sample space\\
\item a probability function, \(P\), which assigns each event in the event space a probability\\
\end{enumerate}
\begin{my_example}
rolling 2 dice cubes:\\
\[
\Omega = \{(m,n) \mid m=1,2,\dots,6, n=1,2,\dots,6\}
\]\\
\end{my_example}
\begin{lemma}
\[
P(\Omega) = 1
\]\\
\end{lemma}
\end{definition}
\section{sum rule}
\label{sec:org2d1bf64}
the \textbf{sum rule} for two random events \(A\) and \(B\) states:\\
\[
P(A\cup B)=P(A)+P(B)-P(A\cap B)
\]\\
if the events are \href{20230127023223-disjoint_events.org}{disjoint}, the formula would simplify to\\
\[
P(A \cup B) = P(A) + P(B)
\]\\
\section{chain rule}
\label{sec:orgcfe57dc}
for two random events \(A\) and \(B\):\\
\[
P(A\cap B) = P(B\mid A) \cdot P(A) = P(A\mid B) \cdot P(B)
\]\\
for more than two events \(A_{1},\dots,A_{n}\) the chain rule extends to\\
\[
P \left(A_{n}\cap \dots \cap A_{1}\right)= P\left(A_{n} \mid A_{n-1}\cap \dots \cap A_{1}\right)\cdot P\left(A_{n-1}\cap \dots \cap A_{1}\right)
\]\\
which by induction may be turned into\\
\[
P\left(A_{n} \cap \dots \cap A_{1}\right)=\prod_{k=1}^{n} P\left(A_{k}\,\Bigg|\,\bigcap_{j=1}^{k-1}A_{j}\right)
\]\\
\section{probability of atleast}
\label{sec:orga8785a9}
\begin{my_example}
\begin{gather*}
  P(\text{at least 3 reds in 5 pulls}) = P(\text{3 reds}) + P(\text{4 reds}) + P(\text{5 reds})\\
  P(\text{at least 1 success}) = 1 - P(\text{all failures})\\
  P(\text{at least 1 failure}) = 1 - P(\text{all successes})
\end{gather*}
\end{my_example}
\end{document}